\documentclass{article}

\usepackage{tabularx}
\usepackage{booktabs}

\title{Problem Statement and Goals\\\progname}

\author{\authname}

\date{}

\input{../Comments.text}
\input{../Common.text}

\begin{document}

\maketitle

\begin{table}[hp]
\caption{Revision History} \label{TblRevisionHistory}
\begin{tabularx}{\textwidth}{llX}
\toprule
\textbf{Date} & \textbf{Developer(s)} & \textbf{Change}\\
\midrule
Date1 & Name(s) & Description of changes\\
Date2 & Name(s) & Description of changes\\
... & ... & ...\\
\bottomrule
\end{tabularx}
\end{table}

\section{Problem Statement}

\wss{You should check your problem statement with the
\href{https://smiths.github.io/capTemplate/Checklists/ProbState-Checklist.pdf}
{problem statement checklist}}. 

\wss{You can change the section headings, as long as you include the required
information.}

\subsection{Problem}

\wss{What problem are you trying to solve?  Why is it important? Do not talk
about how you will solve the problem, only what the problem is.  AI is rarely
part of the problem.  The problem is ``what'' you want to do, not ``how'' you
want to do it.  The problem statement should be written in a way that is
understandable to a non-technical audience.}

Fitness and Finance are two correlated activities, that being they can compound overtime.
However, learning how to deal with finances in a informed manner can be difficult and tedious, there is so many pathways on can take
to manage their money. For example, investing in the stock market...but what makes a stock good to invest in? What are Puts and Calls? What is an ETF, what is a bond, so many questions yet
no single source for an answer. Likewise with fitness, it can be daunting to get started, so many different muscle groups, what is cardio, what is leg day, etc.
It's very important to be both fit and financially literate because much like how good financial and fit decisions can compound over time, bad decisions can too.

\subsection{Inputs and Outputs}

\wss{Characterize the problem in terms of ``high level'' inputs and outputs.  
Use abstraction so that you can avoid details.}
\subsubsection{Inputs}
\item{Biometric sensor data (i.e., steps counted, heart rate, etc)}
\item{Manual user data input (i.e., weight)}
\item{Market real world reference data (i.e., current value of stocks)}


\subsubsection{Outputs}
\item{User market portfolio Value (i.e., current value of owned shares)}
\item{User market portfolio trends and milestones}
\item{User workout trends and milestones}
\item{Market simulation results (i.e., simulated market values given reference input)}
\item{System notifications (i.e., user status from in-app social feeds)}

\subsection{Stakeholders}

\subsection{Environment}

\wss{Hardware and Software Environment for the users.  The developer environment
is summarized as part of the developer plan.}

\section{Goals}
\itme{\textbf{Stock Market Simulation Engine} - Develop an engine which mimicks the real market that users can interact with}
\item{\textbf{A unified interface for biometric data intake} - The system should be designed in such a way that complementing with more biometric data to further refine tracking of user fitness statistics is intuitive to do}

\section{Stretch Goals}

\wss{Goals you hope to get to, but not necessary.  They are also helpful if the
scope of the original project needs to be expanded.}

\section{Custom Documents (CDs)}

\wss{For CAS 741: State whether the project is a research project. This
designation, with the approval (or request) of the instructor, can be modified
over the course of the term.}

\wss{For SE Capstone: List your custom documents (CDs).  Potential CDs include
usability testing, code walkthroughs, user documentation, formal proof,
GenderMag personas, Design Thinking, etc.  (The full list is on the course
outline and in Lecture 02.) Normally the number of CDs will be two (one per
term).  Approval of the CDs will be part of the discussion with the instructor
for approving the project. The CDs, with the approval (or request) of the
instructor, can be modified over the course of the term.}

\newpage{}

\section*{Appendix --- Reflection}

\wss{Not required for CAS 741}

\input{../Reflection.text}

\begin{enumerate}
    \item What went well while writing this deliverable? 
    \item What pain points did you experience during this deliverable, and how
    did you resolve them?
    \item How did you and your team adjust the scope of your goals to ensure
    they are suitable for a Capstone project (not overly ambitious but also of
    appropriate complexity for a senior design project)?
\end{enumerate}  

\end{document}